\section{Charity received and a people renewed}
The two readings meet in the priority of God and the seriousness of fruit. Each begins from gifts the people did not create, faces judgment where those gifts have been misused, and reaches Communion as a source of transformed life. Each refuses to make the neighbor dispensable. Augustine's charity and Chrysostom's one body give the word ``fruit'' a personal and communal measure.

Their differences help locate the difficulty a hearer faces. The first answers the anxious or possessive attempt to turn goodness into a purchase price. Its central sequence is gift, response and nourishment: Christ chooses so that disciples may bear fruit, and the Eucharist stirs charity into action. The second answers despair over a damaged people, or the wish to restore its power without reforming its life. Its sequence is wound, visitation and renewed communion: the same planting is perfected in Christ, and common reception demands common love.

The psalm and Epistle consequently carry different weight. In the first, quickening and thanksgiving explain how response depends on gift. In the second, lament and practiced peace explain how a people asks for and inhabits renewal. The difference is an emphasis, not a quarrel between Augustine and Bellarmine or a claim that Chrysostom taught only one side. Their actual disagreement over the psalm's historical imagery remains a separate question.

Neither Communion option cancels either reading. Lamentations keeps fruitfulness and restoration patient before the Lord whose goodness is sought amid grief. First Corinthians makes their communal consequence explicit: charity is fruit within one body, and restoration is incorporation that obliges the members to love. The prayer after Communion gathers both concerns into the effect of reception upon life. The giver is not finished with his vineyard; the vineyard is not released from bearing his fruit.
